\chapter{Regularization}

\begin{remarkbox}
Loop integrals in QFT are often divergent before they are given a precise
meaning. Regularization is the temporary procedure that makes such expressions
well defined. It does not by itself solve the physical problem; it prepares the
calculation for renormalization.
\end{remarkbox}

\section{Why Divergent Integrals Appear}

Loop diagrams require integration over internal momenta. A typical one-loop
integral has the form
\[
    \int\frac{\dd^4k}{(2\pi)^4}\,
    \frac{1}{k^2-m^2+i\epsilon}.
\]
At large \(|k|\), this behaves too slowly to give a finite answer. The integral
is ultraviolet divergent. Other integrals may diverge at small momentum, giving
infrared divergences.

Divergences arise because field theory has degrees of freedom at arbitrarily
short distances and, for massless particles, arbitrarily long distances. A loop
integral sums over all virtual momenta. Without further definition, this sum may
not exist as an ordinary integral.

Regularization modifies the calculation so that every intermediate expression is
finite or at least well defined. Only then can one isolate divergent pieces and
renormalize them.

\section{Regularization Versus Renormalization}

Regularization and renormalization are related but distinct.

Regularization answers the question:
\[
    \text{How do we make divergent expressions mathematically well defined?}
\]
Renormalization answers the question:
\[
    \text{How do we express predictions in terms of finite measured parameters?}
\]
A regulator is temporary. Physical predictions should not depend on the details
of the regulator after renormalization. If they did, the regulator would be a new
physical assumption rather than a calculational device.

A good regulator should preserve as many symmetries as possible. In gauge
theories, this is especially important because a regulator that violates gauge
invariance can obscure Ward identities and require additional care.

\section{Momentum Cutoff}

The most intuitive regulator is a momentum cutoff. One restricts loop momenta to
satisfy
\[
    |k|<\Lambda.
\]
For example,
\[
    \int\frac{\dd^4k}{(2\pi)^4}\,f(k)
    \quad\longrightarrow\quad
    \int_{|k|<\Lambda}\frac{\dd^4k}{(2\pi)^4}\,f(k).
\]
The cutoff \(\Lambda\) represents a maximum momentum or minimum distance scale.
Divergences then appear as powers or logarithms of \(\Lambda\), such as
\[
    \Lambda^2,
    \qquad
    \ln\Lambda.
\]

The advantage of a cutoff is physical transparency. The disadvantage is that a
simple cutoff can break Lorentz invariance or gauge invariance unless it is
implemented carefully.

\section{Power Divergences and Logarithmic Divergences}

A cutoff makes the degree of divergence visible. An integral may behave as
\[
    I(\Lambda)\sim \Lambda^2
\]
for a quadratic divergence, or
\[
    I(\Lambda)\sim \ln\Lambda
\]
for a logarithmic divergence.

Power divergences are sensitive to the highest momentum scales. Logarithmic
divergences are milder but especially important because they lead directly to
scale dependence and renormalization group running.

For example, a logarithm often appears in the form
\[
    \ln\frac{\Lambda}{\mu},
\]
where \(\mu\) is a reference scale. After renormalization, dependence on \(\mu\)
becomes the scale dependence of renormalized parameters.

\section{Pauli--Villars Regularization}

Pauli--Villars regularization modifies propagators by subtracting fictitious
heavy fields. A scalar propagator may be replaced schematically by
\[
    \frac{i}{k^2-m^2+i\epsilon}
    \quad\longrightarrow\quad
    \frac{i}{k^2-m^2+i\epsilon}
    -\frac{i}{k^2-M^2+i\epsilon},
\]
where \(M\) is a large regulator mass. At large \(k\), the two terms cancel the
leading ultraviolet behavior.

The regulator field is not physical. One takes
\[
    M\to\infty
\]
after renormalization. Pauli--Villars is useful conceptually because it shows
how high-momentum behavior can be softened while preserving Lorentz covariance.

However, Pauli--Villars can be awkward in non-Abelian gauge theories and chiral
theories. It is no longer the most common regulator for modern perturbative
calculations.

\section{Dimensional Regularization}

Dimensional regularization changes the number of spacetime dimensions from four
to
\[
    d=4-\epsilon.
\]
Loop integrals are evaluated as analytic functions of \(d\). Divergences appear
as poles in \(\epsilon\):
\[
    \frac{1}{\epsilon}.
\]
A typical integral becomes
\[
    \int\frac{\dd^4k}{(2\pi)^4}
    \quad\longrightarrow\quad
    \mu^{4-d}\int\frac{\dd^dk}{(2\pi)^d}.
\]
The scale \(\mu\) is introduced so that couplings keep their usual mass
dimensions.

Dimensional regularization is widely used because it preserves Lorentz
invariance and is especially compatible with gauge invariance. It also sets many
power divergences to zero, leaving logarithmic divergences as poles in
\(\epsilon\).

\section{A Standard Dimensionally Regularized Integral}

A common Euclidean-type integral has the form
\[
    \int\frac{\dd^dk}{(2\pi)^d}\,
    \frac{1}{(k^2+\Delta)^\alpha}
    = \frac{1}{(4\pi)^{d/2}}\frac{\Gamma(\alpha-d/2)}{\Gamma(\alpha)}
      \Delta^{d/2-\alpha}.
\]
Divergences appear when the gamma function has a pole. Near \(d=4\), one often
finds terms such as
\[
    \frac{1}{\epsilon}-\gamma_E+\ln(4\pi)-\ln\frac{\Delta}{\mu^2}.
\]
Different subtraction schemes decide which combination of constants is removed.

The most common schemes are minimal subtraction, abbreviated MS, and modified
minimal subtraction, abbreviated \(\overline{\mathrm{MS}}\). In
\(\overline{\mathrm{MS}}\), one subtracts not only \(1/\epsilon\) but also the
standard constants \(-\gamma_E+\ln4\pi\).

\section{Feynman Parameters}

Loop integrals often contain products of denominators. Feynman parameters combine
them. The basic identity is
\[
    \frac{1}{AB}
    = \int_0^1\dd x\,\frac{1}{[xA+(1-x)B]^2}.
\]
More generally,
\[
    \frac{1}{A_1\cdots A_n}
    = (n-1)!\int_0^1\dd x_1\cdots\dd x_n\,
      \delta\left(1-\sum_i x_i\right)
      \frac{1}{(x_1A_1+\cdots+x_nA_n)^n}.
\]
After combining denominators, one shifts the loop momentum and evaluates a
standard Gaussian-type integral.

Feynman parameters are not a regulator by themselves. They are an algebraic tool
that makes regulated loop integrals easier to compute.

\section{Infrared Regularization}

Not all divergences are ultraviolet. In theories with massless particles,
integrals may diverge when loop momentum becomes small. To regulate infrared
divergences, one may introduce a small photon mass \(m_\gamma\), work in
\(d\neq4\) dimensions, or impose an energy resolution for soft radiation.

For example, a soft divergence may appear as
\[
    \ln m_\gamma
\]
if a small photon mass is used. This mass is not physical; it is only an infrared
regulator. Physical inclusive observables must be finite when
\[
    m_\gamma\to0.
\]

Ultraviolet and infrared divergences can appear in the same calculation, but
they have different physical meanings and are removed in different ways.

\section{Regulators and Symmetries}

A regulator can preserve or violate symmetries. This matters. If a regulator
breaks a symmetry that the quantum theory should preserve, additional work is
needed to restore the symmetry in renormalized quantities.

Dimensional regularization is popular because it usually preserves gauge
invariance. A hard momentum cutoff is intuitive but may break gauge invariance if
applied naively. Pauli--Villars can preserve some symmetries but may struggle
with chiral gauge theories.

Sometimes a symmetry cannot be preserved by any regulator. Then the symmetry is
anomalous. Anomalies will be discussed later. For now, the lesson is that the
choice of regulator is not only technical; it can reveal or obscure important
symmetry structure.

\section{Subtraction Schemes}

After regularizing, one subtracts divergent pieces. A subtraction scheme defines
what is removed and what finite parts are included in the definition of
renormalized parameters.

In minimal subtraction, one removes only pole terms such as
\[
    \frac{1}{\epsilon}.
\]
In \(\overline{\mathrm{MS}}\), one removes
\[
    \frac{1}{\epsilon}-\gamma_E+\ln4\pi.
\]
In an on-shell scheme, parameters are fixed by physical pole masses and physical
charges. Different schemes give different numerical values for renormalized
parameters at intermediate stages, but physical predictions agree when computed
consistently.

Scheme dependence is not a contradiction. It reflects the freedom to choose how
finite pieces are assigned between loop corrections and parameter definitions.

\section{What Regularization Does Not Do}

Regularization does not make a theory predictive by itself. A cutoff can make an
integral finite, but the result may depend on the cutoff. Predictivity comes only
after one shows that cutoff dependence can be absorbed into a finite number of
measured parameters, or else interprets the theory as an effective field theory
with cutoff-suppressed corrections.

Regularization also does not mean that infinities are physical. Divergent terms
usually indicate that the parameters in the Lagrangian are not yet related to
physical observables. Renormalization provides that relation.

Thus the correct sequence is
\[
    \text{divergent expression}
    \longrightarrow
    \text{regularized expression}
    \longrightarrow
    \text{renormalized prediction}.
\]

\begin{summarybox}
Regularization is the temporary step that makes divergent loop integrals well
defined. Momentum cutoffs are intuitive but can obscure symmetries.
Pauli--Villars softens high-momentum behavior with heavy regulator fields.
Dimensional regularization evaluates integrals in \(d=4-\epsilon\) dimensions and
represents divergences as poles in \(\epsilon\). Subtraction schemes define which
pieces are absorbed into renormalized parameters. Physical predictions should not
depend on the regulator.
\end{summarybox}

\section{Chapter Checklist}

After reading this chapter, one should be able to:
\begin{itemize}
    \item explain why loop integrals diverge;
    \item distinguish regularization from renormalization;
    \item describe a momentum cutoff;
    \item distinguish power and logarithmic divergences;
    \item explain Pauli--Villars regularization;
    \item explain dimensional regularization and the role of \(d=4-\epsilon\);
    \item use the basic idea of Feynman parameters;
    \item distinguish ultraviolet and infrared regularization;
    \item explain why preserving symmetries matters;
    \item distinguish MS, \(\overline{\mathrm{MS}}\), and on-shell subtraction
    ideas;
    \item state why a regulator is not itself a physical prediction.
\end{itemize}
